\documentclass[10pt,a4paper]{amsart}
\usepackage{amsmath,amssymb,amsthm,graphicx}
\usepackage{verbatim}
\newenvironment{reference}{\comment}{\endcomment}
%\newenvironment{reference}{}{}
\newenvironment{slogan}{\comment}{\endcomment}
\newenvironment{history}{\comment}{\endcomment}

% For commutative diagrams we use Xy-pic
\usepackage[all]{xy}

% We use 2cell for 2-commutative diagrams.
\xyoption{2cell}
% Theorem environments.
%
\theoremstyle{plain}
\newtheorem{theorem}[subsection]{Theorem}
\newtheorem{proposition}[subsection]{Proposition}
\newtheorem{lemma}[subsection]{Lemma}

\theoremstyle{definition}
\newtheorem{definition}[subsection]{Definition}
\newtheorem{example}[subsection]{Example}
\newtheorem{exercise}[subsection]{Exercise}
\newtheorem{situation}[subsection]{Situation}

\theoremstyle{remark}
\newtheorem{remark}[subsection]{Remark}
\newtheorem{remarks}[subsection]{Remarks}

\numberwithin{equation}{subsection}

% Macros
%
\def\lim{\mathop{\mathrm{lim}}\nolimits}
\def\colim{\mathop{\mathrm{colim}}\nolimits}
\def\Spec{\mathop{\mathrm{Spec}}}
\def\Hom{\mathop{\mathrm{Hom}}\nolimits}
\def\Ext{\mathop{\mathrm{Ext}}\nolimits}
\def\SheafHom{\mathop{\mathcal{H}\!\mathit{om}}\nolimits}
\def\SheafExt{\mathop{\mathcal{E}\!\mathit{xt}}\nolimits}
\def\Sch{\mathit{Sch}}
\def\Mor{\mathop{\mathrm{Mor}}\nolimits}
\def\Ob{\mathop{\mathrm{Ob}}\nolimits}
\def\Sh{\mathop{\mathit{Sh}}\nolimits}
\def\NL{\mathop{N\!L}\nolimits}
\def\CH{\mathop{\mathrm{CH}}\nolimits}
\def\proetale{{pro\text{-}\acute{e}tale}}
\def\etale{{\acute{e}tale}}
\def\QCoh{\mathit{QCoh}}
\def\Ker{\mathop{\mathrm{Ker}}}
\def\Im{\mathop{\mathrm{Im}}}
\def\Coker{\mathop{\mathrm{Coker}}}
\def\Coim{\mathop{\mathrm{Coim}}}

% Boxtimes
%
\DeclareMathSymbol{\boxtimes}{\mathbin}{AMSa}{"02}

%
% Macros for moduli stacks/spaces
%
\def\QCohstack{\mathcal{QC}\!\mathit{oh}}
\def\Cohstack{\mathcal{C}\!\mathit{oh}}
\def\Spacesstack{\mathcal{S}\!\mathit{paces}}
\def\Quotfunctor{\mathrm{Quot}}
\def\Hilbfunctor{\mathrm{Hilb}}
\def\Curvesstack{\mathcal{C}\!\mathit{urves}}
\def\Polarizedstack{\mathcal{P}\!\mathit{olarized}}
\def\Complexesstack{\mathcal{C}\!\mathit{omplexes}}
% \Pic is the operator that assigns to X its picard group, usage \Pic(X)
% \Picardstack_{X/B} denotes the Picard stack of X over B
% \Picardfunctor_{X/B} denotes the Picard functor of X over B
\def\Pic{\mathop{\mathrm{Pic}}\nolimits}
\def\Picardstack{\mathcal{P}\!\mathit{ic}}
\def\Picardfunctor{\mathrm{Pic}}
\def\Deformationcategory{\mathcal{D}\!\mathit{ef}}
\usepackage{cleveref}
\crefname{originalsection}{original section}{original sections}
\Crefname{originalsection}{Original section}{Original sections}
\makeatletter
\newlabel{more-algebra-section-injectives-modules}{{01D8}{source}{}{sourcepointer-1}{}}
\newlabel{more-algebra-section-injectives-modules@cref}{{[originalsection][0][]01D8}{[1][0][]source}{}{}{}}
\makeatother
\crefname{originalsection}{original section}{original sections}
\Crefname{originalsection}{Original section}{Original sections}
\makeatletter
\newlabel{sheaves-example-sheaf-product-pointwise}{{006X}{source}{}{sourcepointer-2}{}}
\newlabel{sheaves-example-sheaf-product-pointwise@cref}{{[originalsection][0][]006X}{[1][0][]source}{}{}{}}
\makeatother
\crefname{originallemma}{original lemma}{original lemmas}
\Crefname{originallemma}{Original lemma}{Original lemmas}
\makeatletter
\newlabel{sheaves-lemma-sheaf-subset-stalks}{{0079}{source}{}{sourcepointer-3}{}}
\newlabel{sheaves-lemma-sheaf-subset-stalks@cref}{{[originallemma][0][]0079}{[1][0][]source}{}{}{}}
\makeatother
\crefname{originalsection}{original section}{original sections}
\Crefname{originalsection}{Original section}{Original sections}
\makeatletter
\newlabel{sheaves-section-skyscraper-sheaves}{{0099}{source}{}{sourcepointer-4}{}}
\newlabel{sheaves-section-skyscraper-sheaves@cref}{{[originalsection][0][]0099}{[1][0][]source}{}{}{}}
\makeatother
\crefname{originallemma}{original lemma}{original lemmas}
\Crefname{originallemma}{Original lemma}{Original lemmas}
\makeatletter
\newlabel{sheaves-lemma-stalk-skyscraper-adjoint}{{009C}{source}{}{sourcepointer-5}{}}
\newlabel{sheaves-lemma-stalk-skyscraper-adjoint@cref}{{[originallemma][0][]009C}{[1][0][]source}{}{}{}}
\makeatother
\crefname{originallemma}{original lemma}{original lemmas}
\Crefname{originallemma}{Original lemma}{Original lemmas}
\makeatletter
\newlabel{homology-lemma-product-injectives}{{0137}{source}{}{sourcepointer-6}{}}
\newlabel{homology-lemma-product-injectives@cref}{{[originallemma][0][]0137}{[1][0][]source}{}{}{}}
\makeatother
\title{Abelian sheaves on a topological space have functorial injective embeddings}
\author{The Stacks Project Authors}
\date{stacks-project:injectives.tex:local-snapshot-20261012; GNU FDL 1.2 or later}
\begin{document}
\maketitle
\noindent\textit{Editorial note.} All delivered mathematical text is the author's own, copied from the pinned Stacks Project chapter file sites-cohomology.tex; no statement or proof step was rewritten, reordered or removed.\par\medskip
\noindent\textit{Editorial note.} The proof is the author's complete construction; the injectivity of the target is attributed to a named product lemma carried on this card as a navigation stub, and both asserted properties are otherwise verified in the span.\par\medskip
\noindent\textit{Editorial note.} Every labelled reference occurring inside the counted spans is carried in assembly.external\_references as a non-load-bearing navigation pointer (same-chapter ones at their definition site, cross-chapter ones at the invocation inside the counted span), so no load-bearing mathematics is replaced by a stub. No bibliography citation occurs inside the counted spans.\par\medskip
\noindent\textit{Editorial note.} The same-name check for this card scanned every staged Stacks chapter of this repository -- all sources/stacks-*/pinned-original-source.tex files, the card's own chapter excluded -- and the labels found are recorded in dedup.same\_name\_cross\_chapter on the card.\par\medskip
\noindent\textit{Editorial note.} The preamble is the chapter's own retained minimal preamble (sources/stacks-sites-cohomology/retained-preamble.tex), a verbatim excerpt of the native preamble: the Xy-pic setup and the native theorem declarations with the mathematical macros.\par\medskip
\noindent\textit{Allowed standard prerequisites (editorial).}\begin{itemize}
\item Functorial injective embeddings of abelian groups. Same chapter, the construction of a functorial embedding of an abelian group into an injective group. Supplies the embeddings of the groups that appear in the skyscraper construction.
\item Skyscraper sheaves and their adjunction with the stalk functor. Sheaves, Lemma sheaves-lemma-stalk-skyscraper-adjoint. Supplies the adjunction that identifies the sheaf constructed from the injective groups at the points as a product of skyscraper sheaves.
\item Products of injective objects. Homology, Lemma homology-lemma-product-injectives. Supplies that a product of injective objects is injective, which gives injectivity of the constructed sheaf.
\end{itemize}\medskip
\ifdefined\hypertarget\hypertarget{sourcepointer-1}{}\fi
\noindent\textit{Source reference.} Original section 01D8: sources/stacks-more-algebra/pinned-original-source.tex, source lines 13617--13617. Navigation pointer invoked at proof line 623; the counted span contains the pointer only and the cited result is not reproduced inside it, so the pointer is registered as a non-load-bearing navigation stub in the established form for Stacks material. No bibliography citation is involved.\par\medskip
\ifdefined\hypertarget\hypertarget{sourcepointer-2}{}\fi
\noindent\textit{Source reference.} Original section 006X: sources/stacks-sheaves/pinned-original-source.tex, source lines 814--814. Navigation pointer invoked at proof line 667; the counted span contains the pointer only and the cited result is not reproduced inside it, so the pointer is registered as a non-load-bearing navigation stub in the established form for Stacks material. No bibliography citation is involved.\par\medskip
\ifdefined\hypertarget\hypertarget{sourcepointer-3}{}\fi
\noindent\textit{Source reference.} Original lemma 0079: sources/stacks-sheaves/pinned-original-source.tex, source lines 1163--1163. Navigation pointer invoked at proof line 677; the counted span contains the pointer only and the cited result is not reproduced inside it, so the pointer is registered as a non-load-bearing navigation stub in the established form for Stacks material. No bibliography citation is involved.\par\medskip
\ifdefined\hypertarget\hypertarget{sourcepointer-4}{}\fi
\noindent\textit{Source reference.} Original section 0099: sources/stacks-sheaves/pinned-original-source.tex, source lines 3487--3487. Navigation pointer invoked at proof line 689; the counted span contains the pointer only and the cited result is not reproduced inside it, so the pointer is registered as a non-load-bearing navigation stub in the established form for Stacks material. No bibliography citation is involved.\par\medskip
\ifdefined\hypertarget\hypertarget{sourcepointer-5}{}\fi
\noindent\textit{Source reference.} Original lemma 009C: sources/stacks-sheaves/pinned-original-source.tex, source lines 3549--3549. Navigation pointer invoked at proof line 696; the counted span contains the pointer only and the cited result is not reproduced inside it, so the pointer is registered as a non-load-bearing navigation stub in the established form for Stacks material. No bibliography citation is involved.\par\medskip
\ifdefined\hypertarget\hypertarget{sourcepointer-6}{}\fi
\noindent\textit{Source reference.} Original lemma 0137: sources/stacks-homology/pinned-original-source.tex, source lines 7335--7335. Navigation pointer invoked at proof line 699; the counted span contains the pointer only and the cited result is not reproduced inside it, so the pointer is registered as a non-load-bearing navigation stub in the established form for Stacks material. No bibliography citation is involved.\par\medskip
\begin{lemma}
\label{lemma-abelian-sheaves-space}
Let $X$ be a topological space.
The category of abelian sheaves on $X$ has enough injectives.
In fact it has functorial injective embeddings.
\end{lemma}

\begin{proof}
For an abelian group $A$ we denote $j : A \to J(A)$ the functorial
injective embedding constructed in
More on Algebra, Section \ref{more-algebra-section-injectives-modules}.
Let $\mathcal{F}$ be an abelian sheaf on $X$.
By Sheaves, Example \ref{sheaves-example-sheaf-product-pointwise}
the assignment
$$
\mathcal{I} : U \mapsto
\mathcal{I}(U) = \prod\nolimits_{x\in U} J(\mathcal{F}_x)
$$
is an abelian sheaf. There is a canonical map $\mathcal{F} \to \mathcal{I}$
given by mapping $s \in \mathcal{F}(U)$ to $\prod_{x \in U} j(s_x)$
where $s_x \in \mathcal{F}_x$ denotes the germ of $s$ at $x$.
This map is injective, see
Sheaves, Lemma \ref{sheaves-lemma-sheaf-subset-stalks}
for example.

\medskip\noindent
It remains to prove the following: Given a rule
$x \mapsto I_x$ which assigns to each point $x \in X$ an injective
abelian group the sheaf $\mathcal{I} : U \mapsto \prod_{x \in U} I_x$
is injective. Note that
$$
\mathcal{I} = \prod\nolimits_{x \in X} i_{x, *}I_x
$$
is the product of the skyscraper sheaves $i_{x, *}I_x$ (see
Sheaves, Section \ref{sheaves-section-skyscraper-sheaves} for notation.)
We have
$$
\Mor_{\textit{Ab}}(\mathcal{F}_x, I_x)
=
\Mor_{\textit{Ab}(X)}(\mathcal{F}, i_{x, *}I_x).
$$
see Sheaves, Lemma \ref{sheaves-lemma-stalk-skyscraper-adjoint}. Hence it is
clear that each $i_{x, *}I_x$ is injective. Hence the injectivity of
$\mathcal{I}$ follows from
Homology, Lemma \ref{homology-lemma-product-injectives}.
\end{proof}

\end{document}
